\documentclass[10pt,a4paper]{article}
\usepackage[a4paper,margin=2.1cm]{geometry}
\usepackage[T1]{fontenc}
\usepackage[utf8]{inputenc}
\usepackage{lmodern}
\usepackage{microtype}
\usepackage{xcolor}
\usepackage{enumitem}
\usepackage{tabularx}
\usepackage{longtable}
\usepackage{array}
\usepackage{amssymb}
\usepackage{hyperref}
\usepackage{fancyhdr}
\definecolor{tudblue}{RGB}{0,90,140}
\definecolor{draftred}{RGB}{165,35,35}
\definecolor{lightblue}{RGB}{232,242,247}
\hypersetup{colorlinks=true,linkcolor=tudblue,urlcolor=tudblue}
\setlength{\parindent}{0pt}
\setlength{\parskip}{0.55em}
\setlist{itemsep=0.25em,topsep=0.35em}
\renewcommand{\arraystretch}{1.22}
\pagestyle{fancy}
\fancyhf{}
\lhead{TU Darmstadt Ethics Commission}
\rhead{Draft for supervisor review}
\cfoot{\thepage}
\newcommand{\checked}{\ensuremath{\boxtimes}}
\newcommand{\unchecked}{\ensuremath{\square}}
\newcommand{\confirm}[1]{\textcolor{draftred}{\textbf{TO CONFIRM:} #1}}
\newcommand{\field}[2]{\textbf{#1}\par #2\par}
\newcommand{\draftbanner}{%
  \begin{center}
  \fcolorbox{draftred}{white}{\parbox{0.91\textwidth}{\centering
  \textcolor{draftred}{\textbf{DRAFT -- NOT FOR SUBMISSION OR DATA COLLECTION}}\\
  Administrative and data-management fields marked ``TO CONFIRM'' require supervisor approval.}}
  \end{center}}
\newcommand{\sectionbox}[1]{%
  \vspace{0.4em}\noindent\colorbox{lightblue}{\parbox{\dimexpr\linewidth-2\fboxsep\relax}{\textbf{#1}}}\vspace{0.35em}}

\begin{document}

\begin{center}
  {\Large\bfseries Ethics Application 2026 -- Completed Draft}\\[0.3em]
  {\large Local Trait Inference on Smartphones}
\end{center}
\draftbanner

\section*{Brief Description of the Research Project}

\sectionbox{Principal Investigator}
\field{Academic title and name}{\confirm{Principal investigator.}}
\field{Research group}{Science and Technology for Peace and Security (PEASEC); collaboration with Work and Engineering Psychology to be confirmed.}
\field{Department}{Department of Computer Science, Technical University of Darmstadt; possible involvement of the Department of Human Sciences.}

\sectionbox{Information on the Research Project}
\field{Working title}{Local Trait Inference on Smartphones: User Perceptions of On-Device AI Profiling}
\field{Brief description (maximum 700 characters)}{This bachelor-thesis study examines how adults perceive privacy-relevant profiles generated by an AI model directly on their Android phone. Participants select approximately 20 of their own images or PDF documents. Raw files remain on the device; the app displays probabilistic file-level inferences and a combined profile. Participants complete privacy, technology-affinity, demographic, per-inference, and final perception questions. Optional front-camera reaction tracking stores emotion labels but no images. The study evaluates expected capability, perceived accuracy, sensitivity, surprise, concern, and willingness to grant file access.}
\field{Planned project running time}{26 May 2026 to 30 November 2026.}
\field{Estimated timeframe for data collection}{1 September 2026 to 31 October 2026 at the latest, and only after a positive ethics vote and final approval of all study materials.}
\field{Other participating researchers}{Atabey Tarik Akkum (student researcher, Bachelor of Science in Computer Science); Simon Althaus, M.Sc. (thesis supervisor, PEASEC); Dr. rer. nat. Nina Gerber (psychological and user-study consultation; formal role to be confirmed); \confirm{principal investigator and any additional data analysts}.}

\field{The study is part of}{\checked\ B.A. Thesis \quad \unchecked\ M.A. Thesis \quad \unchecked\ Dissertation \quad \unchecked\ Course \quad \unchecked\ Other}
\field{Motivation for the application}{\checked\ The project touches on ethically relevant issues, specifically personal data, privacy-sensitive AI inferences, optional camera-based emotion classification, and the possible incidental presence of third-party information in participant-selected files.}

\newpage
\section*{Checklist}
\begin{tabularx}{\textwidth}{>{\raggedright\arraybackslash}Xcc}
\textbf{Question} & \textbf{Yes} & \textbf{No}\\\hline
Physical stress or hazards? & \unchecked & \checked\\
Pain or mental stress through strong emotions beyond everyday life? & \unchecked & \checked\\
Deception or manipulation of participants? & \unchecked & \checked\\
Patients included? & \unchecked & \checked\\
Minors included? & \unchecked & \checked\\
Possible exposure of individuals, minorities, or vulnerable groups through statistical linking? & \unchecked & \checked\\
Might third parties be involved or affected? & \checked & \unchecked\\
Is personal data collected from participants? & \checked & \unchecked\\
Hazardous substances used or produced? & \unchecked & \checked\\
Publication without guaranteed participant anonymity? & \unchecked & \checked\\
Research with commercially created datasets? & \unchecked & \checked\\
Uncertainty about data origin or lawful collection? & \unchecked & \checked\\
Commercial project partner also using collected data? & \unchecked & \checked\\
In the applicants' opinion, is the study ethically uncritical? & \unchecked & \checked\\
\end{tabularx}

\textbf{Checklist rationale.} The study does not intentionally induce distress beyond ordinary encounters with privacy notices or personalized digital services. Some participants may nevertheless experience concern, surprise, or mild discomfort after seeing probabilistic personal inferences. Personal data are collected. Third-party information could appear in selected files despite instructions; therefore this issue is conservatively marked ``Yes.'' The proposal is not classified as ethically uncritical because it deliberately studies privacy-sensitive profiling and includes optional camera-based processing.

\newpage
\section*{Application: Detailed Description of the Research Project}

\subsection*{Aim and research questions}
Modern mobile devices can run multimodal AI locally. This allows an application to infer interests, activities, plans, education or occupational context, and other personal signals without transmitting the underlying images or documents. The study investigates:
\begin{enumerate}
  \item what privacy-relevant information a local multimodal model can infer from participant-selected smartphone files;
  \item how participants rate individual inferences in terms of unexpectedness, sensitivity, and perceived accuracy;
  \item whether experiencing local profiling changes expectations about smartphone AI, concern about repeated background analysis, reassurance associated with the statement ``data stays on the device,'' and willingness to grant broad file access; and
  \item how privacy concern and affinity for technology interaction relate to these perceptions.
\end{enumerate}

The optional camera module provides exploratory emotion-label observations during study interaction. It is not used to identify participants or diagnose psychological or medical conditions.

\subsection*{Design and participant tasks}
The study uses a within-participant app-based design. Approximately 100 adults will use the Android research prototype on an Android smartphone. The planned procedure is:
\begin{enumerate}
  \item Read the participant information and provide informed consent through unticked in-app checkboxes. Participation cannot proceed without the required consent. Optional front-camera tracking has a separate opt-in and is not required.
  \item Complete demographic questions, IUIPC-8, the nine-item Affinity for Technology Interaction (ATI) scale, and four baseline questions about local AI expectations.
  \item Select approximately 20 personally owned image files or PDF documents. The app instructs participants not to select files depicting identifiable third parties or containing third-party personal information. Files can be skipped and the study can be stopped at any time.
  \item Analyze the files locally with Gemma 4 E2B. Images are analyzed directly. For PDFs, only a rendered preview of the first page is analyzed. The model produces structured probabilistic inferences; participants are explicitly told that these are guesses rather than confirmed facts.
  \item After each successful inference, rate how unexpected, sensitive, and accurate it appears on five-point verbal scales. Completed inferences can be evaluated while later files are still being processed.
  \item Generate a compact combined profile from the file-level outputs and answer six final questions about local AI capability, repeated background profiling, data-locality reassurance, undisclosed summary transmission, commercial use, and future file-access decisions.
  \item Export one pseudonymized study record containing questionnaire responses, model outputs, ratings, technical metadata, optional emotion labels, and the combined profile. Raw selected files and camera images are not included.
\end{enumerate}

\subsection*{Participants, recruitment, setting, duration, and compensation}
Inclusion criteria are age 18 years or older, capacity to provide informed consent, and ability to use the Android prototype. Minors, patients, and vulnerable groups are not specifically recruited. No medical or psychological condition is assessed.

\confirm{Recruitment route, study setting, and compensation. Current proposal: convenience sampling among adults accessible to the research team, without coercion and without recruitment from a direct dependency relationship. Confirm whether TU channels or participant pools will also be used and whether compensation or course credit is offered.}

The active questionnaire and analysis procedure is expected to require approximately 25--40 minutes. A first-time model download of approximately 2.5 GB may add waiting time depending on the network and can run while baseline questions are completed. Wi-Fi is recommended. The participant-facing duration will be finalized after a short pilot.

\subsection*{Potential benefits}
Participants receive no guaranteed personal benefit. The study may increase understanding of current on-device AI capabilities. Scientifically, it contributes evidence about how local processing affects privacy expectations and whether avoiding raw-file upload is sufficient to reassure users when derived profiles can still be transmitted.

\subsection*{Risks and burden}
There are no physical risks. The principal foreseeable burden is mild concern, surprise, embarrassment, or discomfort caused by inaccurate or sensitive-seeming inferences. Mitigations are voluntary file selection, explicit uncertainty language, the ability to skip files or stop immediately, no requirement to select highly sensitive material, and contact details for questions or complaints. No deception is used. The true purpose of the study is disclosed before participation.

Participants are warned not to include images of other identifiable people or documents containing third-party personal information. If such a file is selected accidentally, it can be removed and will not be exported as a raw file. The final implementation will replace filenames and local URIs with random file identifiers before export.

\subsection*{Optional front-camera reaction tracking}
Participants may separately enable front-camera emotion tracking. If enabled, a low-resolution frame is sampled approximately every 12 seconds while the relevant study interface is active. Each frame is analyzed locally by a task-specific emotion-classification model. Only timestamp, study event, model identifier, predicted emotion label, and confidence are retained. No audio is recorded. Frame bytes and temporary camera files will be deleted immediately after each classification and will not be exported. Declining or stopping this feature has no effect on participation.

\section*{Handling of Checklist Issues Marked ``Yes''}
\textbf{Personal data.} Questionnaire responses, consent decisions, per-inference ratings, model-generated inferences, a combined profile, technical metadata, and optional emotion labels are personal data. Data minimization, pseudonymization, restricted access, purpose limitation, and defined deletion periods will be used. Raw selected files remain local and are not exported.

\textbf{Potentially affected third parties.} A participant-selected file could contain another person's image or information. Instructions prohibit such selection unless the third party has independently consented. Participants can remove files before analysis. No raw file is exported. Random file identifiers prevent filenames from disclosing third-party information.

\textbf{Ethical sensitivity.} The study deliberately makes privacy consequences visible. It avoids deception, coercion, diagnostic claims, and intentionally alarming language. Outputs are labelled as probabilistic. Participants can terminate without disadvantage.

\section*{Regulations and Standards}
The study follows the EU General Data Protection Regulation (GDPR), the Hessian Data Protection and Freedom of Information Act (HDSIG), informed-consent and data-minimization principles, the DFG Guidelines for Safeguarding Good Research Practice, and applicable TU Darmstadt research-data and information-security requirements. The study will begin only after ethics approval. The final data flow and storage arrangement will be reviewed with the TU Darmstadt Data Protection Officer if required.

\newpage
\section*{Data Protection}

\field{Methods and data collection instruments}{App-based experiment and questionnaires using an Android research prototype. Instruments comprise in-app consent, demographic items, IUIPC-8, ATI, custom baseline and final questions, per-inference verbal rating scales, local Gemma inference, technical runtime metadata, and optional local front-camera emotion classification.}

\field{Affected groups and number}{Approximately 100 adults aged 18 or older. No targeted recruitment of minors, patients, or vulnerable groups. Third parties are potentially affected only if a participant disregards the instruction not to select files containing other identifiable people or their personal data.}

\field{How consent is obtained}{Before any study data are collected, participants receive the full information sheet in the app and confirm required statements using initially unticked checkboxes. They consent to participation, local processing of selected files, and pseudonymized export of questionnaire and inference data. Camera tracking requires an additional optional opt-in. Consent and participation can be withdrawn without disadvantage.}

\field{Types of personal data}{Random study identifier; age group, gender including ``prefer not to say,'' and education category; IUIPC-8 and ATI responses; local-AI expectation items; consent decisions; per-inference unexpectedness, sensitivity, and accuracy ratings; structured model inferences and combined profile; model/runtime metadata and latency; optional timestamped emotion label and confidence. The final export will not contain raw files, camera frames, original filenames, local file paths, names, email addresses, or phone numbers.}

\field{Deletion periods}{a) Raw selected files are processed locally and are never copied to the research server. Rendered PDF previews and camera frames are deleted immediately after use or at latest when the local study session ends. b) \confirm{retention period for pseudonymized and anonymized research data; proposed distinction between working data and the data underlying thesis/publications}. c) \confirm{retention and separate storage period for consent records}. d) No audio is collected; temporary camera images are deleted immediately, and only derived labels are retained under the pseudonymous study ID.}

\field{Pseudonymization/anonymization}{The app assigns a random study ID unrelated to the participant's identity. Exported files use randomized per-file identifiers. Recruitment contact details, if any, are stored separately and are never joined with study records. Analyses and publications use aggregated results. A participant can use the study ID to request deletion until the agreed anonymization cutoff.}

\field{Local or online analysis and security}{Raw-file and camera-frame analysis occurs locally on the Android phone. One JSON study record is transmitted only after consent using encrypted transport to \confirm{the approved TU/PEASEC server and technical safeguards}. Stored records are pseudonymized, protected by role-based access, and included in institutional backup and deletion procedures. The production study will avoid third-party model downloads that disclose participant IP addresses, for example through TU-hosted distribution or supervised preinstallation.}

\field{Persons analyzing the data and confidentiality}{Atabey Tarik Akkum and the confirmed thesis supervisors/researchers will analyze the data. \confirm{final named access list}. Access is limited to the study team. All researchers are bound by institutional confidentiality and data-protection obligations.}

\field{Deadline for deletion requests}{Participants can request deletion using their study ID until \confirm{the date or processing milestone after which records are irreversibly anonymized and can no longer be linked to an individual participation}.}

\field{Subsequent use}{No subsequent use or publication of raw participant files or camera images is planned. Anonymized aggregate results may be used in the bachelor thesis and related scientific publications/presentations. \confirm{whether pseudonymized record-level data may be retained for follow-up research; default recommendation: no reuse beyond the approved project without renewed legal and ethical review}.}

\section*{Attachments}
\checked\ Informed consent and participant information\\
\checked\ Questionnaires\\
\unchecked\ Research proposal\\
\unchecked\ Images/screenshots\\
\unchecked\ Other

\section*{Signature}
\confirm{The confirmed principal investigator must review and sign.}

\vspace{2.5em}
\rule{0.8\textwidth}{0.4pt}\\
Date and signature of the Principal Investigator

\end{document}
